% !TEX program = lualatex
% Transparencias de Fundamentos de las Matemáticas II
% Clase 14: Números reales II

\documentclass[aspectratio=169,11pt]{beamer}

\usepackage{fontspec}
\usepackage[spanish,es-nodecimaldot]{babel}
\usepackage{amsmath,amssymb,mathtools}
\usepackage{booktabs}
\usepackage{csquotes}
\usepackage{xcolor}

\usetheme{Madrid}
\usecolortheme{default}
\setbeamertemplate{navigation symbols}{}
\setbeamertemplate{footline}[frame number]

\definecolor{FdMBlue}{RGB}{20,55,100}
\definecolor{FdMRed}{RGB}{130,30,30}
\definecolor{FdMGreen}{RGB}{25,100,60}

\setbeamercolor{structure}{fg=FdMBlue}
\setbeamercolor{title}{fg=white,bg=FdMBlue}
\setbeamercolor{frametitle}{fg=white,bg=FdMBlue}
\setbeamercolor{block title}{fg=white,bg=FdMBlue}
\setbeamercolor{block body}{bg=blue!3}
\setbeamercolor{alerted text}{fg=FdMRed}

\setmainfont{Latin Modern Roman}
\setsansfont{Latin Modern Sans}
\setmonofont{Latin Modern Mono}

\newcommand{\N}{\mathbb{N}}
\newcommand{\Z}{\mathbb{Z}}
\newcommand{\Q}{\mathbb{Q}}
\newcommand{\R}{\mathbb{R}}
\newcommand{\C}{\mathbb{C}}
\newcommand{\Pow}{\mathcal{P}}
\newcommand{\id}{\operatorname{id}}
\newcommand{\mcd}{\operatorname{mcd}}
\newcommand{\mcm}{\operatorname{mcm}}
\newcommand{\im}{\operatorname{Im}}

\title[FdM II -- Clase 14]{Números reales II}
\subtitle{Valor absoluto, distancia, supremo e ínfimo}
\author{Antonio Falcó Montesinos}
\institute{Universidad CEU Cardenal Herrera}
\date{Fundamentos de las Matemáticas II}

\begin{document}

\begin{frame}
  \titlepage
\end{frame}

\begin{frame}{Objetivos de la clase}
\begin{itemize}
  \item Definir valor absoluto y distancia en \(\mathbb{R}\).
  \item Comprender cotas, supremo e ínfimo.
  \item Distinguir máximo y supremo.
\end{itemize}
\end{frame}

\begin{frame}{Mapa de la clase}
\tableofcontents
\end{frame}

\section{Valor absoluto y distancia}

\begin{frame}{Valor absoluto y distancia}
\begin{block}{Valor absoluto}
\begin{itemize}
  \item \(|x|=x\) si \(x\geq 0\).
  \item \(|x|=-x\) si \(x<0\).
  \item Mide la distancia de \(x\) al cero.
\end{itemize}
\end{block}
\end{frame}

\begin{frame}{Valor absoluto y distancia}
\begin{block}{Distancia}
\begin{itemize}
  \item La distancia entre \(x\) e \(y\) se define como \(|x-y|\).
  \item Siempre es no negativa.
  \item Es cero si y solo si \(x=y\).
\end{itemize}
\end{block}
\end{frame}

\begin{frame}{Valor absoluto y distancia}
\begin{block}{Desigualdad triangular}
\begin{itemize}
  \item Para todo \(x,y\in\mathbb{R}\), \(|x+y|\leq |x|+|y|\).
  \item Expresa que el camino directo no supera la suma de caminos parciales.
  \item Es una herramienta básica en análisis.
\end{itemize}
\end{block}
\end{frame}

\section{Supremo e ínfimo}

\begin{frame}{Supremo e ínfimo}
\begin{block}{Cotas}
\begin{itemize}
  \item Una cota superior de \(A\) es un número \(M\) tal que \(a\leq M\) para todo \(a\in A\).
  \item Una cota inferior se define de forma análoga.
  \item Un conjunto puede tener muchas cotas.
\end{itemize}
\end{block}
\end{frame}

\begin{frame}{Supremo e ínfimo}
\begin{block}{Supremo}
\begin{itemize}
  \item El supremo es la menor cota superior.
  \item Puede no pertenecer al conjunto.
  \item Si pertenece al conjunto, coincide con el máximo.
\end{itemize}
\end{block}
\end{frame}

\begin{frame}{Supremo e ínfimo}
\begin{block}{Ínfimo}
\begin{itemize}
  \item El ínfimo es la mayor cota inferior.
  \item Puede no pertenecer al conjunto.
  \item Si pertenece al conjunto, coincide con el mínimo.
\end{itemize}
\end{block}
\end{frame}

\section{Profundización}

\begin{frame}{Profundización}
\begin{block}{Ejemplo}
\begin{itemize}
  \item Para \(A=(0,1)\), \(\sup A=1\) e \(\inf A=0\).
  \item Ni \(0\) ni \(1\) pertenecen a \(A\).
  \item Por tanto, \(A\) no tiene máximo ni mínimo.
\end{itemize}
\end{block}
\end{frame}

\begin{frame}{Profundización}
\begin{block}{Ejemplo cerrado}
\begin{itemize}
  \item Para \(B=[0,1]\), \(\sup B=1\) y \(\inf B=0\).
  \item Aquí ambos pertenecen a \(B\).
  \item Por tanto, \(\max B=1\) y \(\min B=0\).
\end{itemize}
\end{block}
\end{frame}

\begin{frame}{Profundización}
\begin{block}{Errores frecuentes}
\begin{itemize}
  \item Confundir máximo con supremo.
  \item Pensar que todo conjunto acotado tiene máximo.
  \item Tratar \(+\infty\) como si fuera un real.
\end{itemize}
\end{block}
\end{frame}


\section{Detalle y práctica}

\begin{frame}{Detalle y práctica}
\begin{block}{Supremo e ínfimo}
\begin{itemize}
  \item Una cota superior de \(A\subset\mathbb{R}\) es un número \(M\) tal que \(a\leq M\) para todo \(a\in A\).
  \item El supremo es la menor de las cotas superiores.
  \item Puede pertenecer o no pertenecer al conjunto.
\end{itemize}
\end{block}
\end{frame}

\begin{frame}{Detalle y práctica}
\begin{block}{Actividad breve}
\begin{itemize}
  \item Hallar el supremo y el ínfimo de \((0,1)\).
  \item Hallar el máximo y el mínimo de \([0,1]\).
  \item Comparar qué cambia entre intervalo abierto e intervalo cerrado.
\end{itemize}
\end{block}
\end{frame}

\begin{frame}{Cierre}
\begin{itemize}
  \item El valor absoluto mide distancia.
  \item Supremo e ínfimo formalizan cotas óptimas.
  \item Máximo y mínimo exigen pertenencia al conjunto.
\end{itemize}
\end{frame}

\begin{frame}{Trabajo recomendado}
\begin{itemize}
  \item Releer las secciones correspondientes del manual.
  \item Identificar definiciones, símbolos nuevos y enunciados principales.
  \item Preparar dudas y ejemplos para el seminario asociado.
\end{itemize}
\end{frame}

\end{document}
